\subsubsection{Identificación de Recursos}
% Guia: recursos humanos, de hardware, software, licencias e infraestructura
% requeridos.
El desarrollo del proyecto cuenta con los siguientes recursos:

\textbf{Recursos de software.} El proyecto se apoya en herramientas de software de código abierto, por lo que no se requiere de la adquisición de alguna licencia de pago. Entre las herramientas utilizadas se encuentran:

\begin{itemize}
    \item Lenguaje de programacion Python, junto con las librerías Flask (framework web), pandas, scikit-learn, XGBoost, statsmodels y SQLAlchemy, entre otras, para el desarrollo del sistema y del pipeline de predicción.
    \item MySQL como motor de base de datos relacional, y MySQL Workbench como herramienta de administración y consulta.
    \item Git y GitHub para el control de versiones del código fuente, junto con GitHub Actions para la ejecución automatizada de pruebas (integración continua).
    \item \textit{pytest} para la implementación de las pruebas automatizadas del sistema.
    \item Visual Studio Code como editor de código para el desarrollo del proyecto.
\end{itemize}


\textbf{Recursos humanos.} En cuanto a recursos humanos, el proyecto fue desarrollado individualmente por una única persona que tomó los roles de investigación, analista de datos y desarrolladora del sistema, esto junto a la orientacion del asesor del proyecto semanalmente.


\textbf{Recursos de hardware.} Tanto el desarrollo del sistema como las pruebas se realizaron en un computador portatil con procesador Intel Core i5-6200U (2.30GHz), 16GB de memoria RAM y almacenamiento de 238GB en unidad de estado solido (SSD), sin requerir hardware especializado (como unidades de procesamiento gráfico dedicadas) para el entrenamiento de los modelos de predicción, dado el volumen de datos manejado en el caso de estudio (856 productos, con historial mensual limitado al año 2021).


\textbf{Infraestructura y despliegue.} El despliegue del sistema se encuentra en la plataforma de hosting en la nube Railway, bajo un plan de uso gratuito, el cual es suficiente para el alcance del proyecto.


\subsubsection{Análisis de Viabilidad}
% Guia: viabilidad tecnica, economica, operativa y legal de la solucion propuesta.
El análisis de viabilidad del proyecto se realiza teniendo en cuenta cuatro aspectos: técnica, económica, operativa y legal. A continuación se describen cada uno de ellos:

\textbf{Viabilidad técnica.} El proyecto se considera tecnicamente viable, dado que se apoya únicamente en tecnologias de código abierto que se encuentran ampliamente documentadas.

\textbf{Viabilidad económica.} Debido a los límites con los que cuenta el proyecto, este se desarrollo implementando herramientas de acceso libre, que sean gratuitas o al menos tengan un plan gratuito, de manera que el costo económico del sistema es mínimo.

\textbf{Viabilidad operativa.} El diseño del sistema cuenta con una aplicación web que es accesible desde un navegador, no es necesario que el usuario final tenga que instalar un software adicional, por lo que es fácil de usar y adoptar en la operación diaria. Además, la interfaz de usuario cuenta con un diseño intiutivo y fácil de usar para cualquier tipo de usuario.

\textbf{Viabilidad legal.} Los datos históricos reales que fueron usados para probar el sistema fueron proporcionados con fines exclusivamente académicos. El sistema cuenta con mecanismos de seguridad que garantizan la confidencialidad de la información manejada.